% ============================================================================ 1 · Descripción general
\section{Descripción general del desafío}
\label{sec:desafio}

\textbf{Título.} Ford pide una metodología de IA para predecir eventos de degradación de la eficiencia de combustión e
ingreso de oxígeno al motor, a partir de datos de vehículos conectados.

\textbf{Qué entendimos.} En un diésel, el componente que se degrada por el uso y afecta a la vez la combustión y el flujo
de gases es el filtro de partículas (DPF). El filtro atrapa el hollín y lo quema para limpiarse (regenera), pero para eso
necesita temperatura: un tramo de ruta con el motor caliente. Con viajes cortos, la regeneración empieza y no termina, el
hollín se acumula y, si sigue así, el auto entra en modo de protección y termina en el taller. Para el cliente, eso es un
testigo que nadie le explicó y la sensación de que nadie le avisó.

\textbf{Qué encontramos en los datos.}
\begin{itemize}
  \item \textbf{El evento es del filtro, con otro nombre.} \emph{TripSummary} no trae las columnas del DPF del anexo, pero sí
    \texttt{AirRegenerationStart/\allowbreak{}End} y \texttt{AirFilterStart/\allowbreak{}End}, que el anexo no menciona. Son la foto, al inicio y al
    final de cada viaje, de la acumulación (0–100) y de los nueve niveles de mensaje del filtro que trae
    \emph{DynamicInformation} (coinciden en el 99,8 y el 99,97\,\% de los pares). Ford no lo confirmó, así que hablamos de
    “nivel de acumulación del filtro”.
  \item \textbf{Es saturación por uso, no desgaste.} El evento cae a una mediana de \textasciitilde12.800 km, y la ceniza
    es un fenómeno de más de 100.000 km. Buscamos una carga irreversible medible (la premisa de un proceso de degradación
    tipo gamma) y no la hay. Eso orientó todo hacia los hábitos de manejo y el estado reciente del filtro.
  \item \textbf{No todo “uso indebido” se sostiene.} Comparando contra sanos \emph{del mismo mercado}, lo que crece hacia el
    evento es estar detenido con el motor en marcha y no llegar a temperatura. La velocidad y la fracción de viajes
    cortos no distinguen dentro del mercado: con la primera entrega, la velocidad parecía la mejor señal y era Colombia
    contra Chile.
\end{itemize}

\textbf{Qué hicimos.} Un modelo de deep learning que lee la secuencia de los últimos 1.000 km del auto y un producto que
convierte cada alerta en una acción. Al conductor no le hablamos de “uso indebido” ni de fallas: le hablamos de hábitos
que puede cambiar, comparados con autos sanos de su zona.
